% P6 manuscript — target: Journal of Computer-Aided Molecular Design (JCAMD) (Springer Nature)
\documentclass[11pt,a4paper]{article}

\usepackage[margin=1in]{geometry}
\usepackage[T1]{fontenc}
\usepackage[utf8]{inputenc}
\usepackage{lmodern}
\usepackage{amsmath,amssymb}

\usepackage{booktabs}
\usepackage{multirow}
\usepackage{graphicx}
\graphicspath{{../results/figures/}}
\usepackage[colorlinks=true,linkcolor=blue,citecolor=blue,urlcolor=blue]{hyperref}
\usepackage{microtype}
\usepackage[numbers,sort&compress]{natbib}
\usepackage{xcolor}
\usepackage{siunitx}
\sisetup{
 detect-all          = true,
 group-minimum-digits = 4,
 range-phrase        = {--},
 range-units         = single,
 separate-uncertainty = true,
 retain-explicit-plus = true,
}
\usepackage[nameinlink,noabbrev]{cleveref}
\crefname{table}{Table}{Tables}
\Crefname{table}{Table}{Tables}
\crefname{figure}{Figure}{Figures}
\Crefname{figure}{Figure}{Figures}
\crefname{equation}{Eq.}{Eqs.}
\crefname{section}{Section}{Sections}
\crefname{subsection}{Section}{Sections}

\title{What Transfers Across Chemical Boundaries? Leakage-Aware Structure--Phenotype Prediction of Observed MoA-Associated Labels}



\date{}

\begin{document}

{\centering\Large\bfseries What Transfers Across Chemical Boundaries? Leakage-Aware Structure--Phenotype Prediction of Observed MoA-Associated Labels\par}
\vspace{1em}
{\centering
Myke Vital Sao Temgoua\thanks{Corresponding author: \texttt{myke-vital.sao@facsciences-uy1.cm}}$^{1}$,
Jean-Pierre Tchapet Njafa$^{1}$,
Penabei Samafou$^{2}$,
Wilfred Fon Mbacham$^{3}$,
Serge Guy Nana Engo$^{1}$\\[0.5em]
{\small $^{1}$Department of Physics, Faculty of Science, University of Yaound\'e I, P.O. Box 812, Yaound\'e, Cameroon}\\
{\small $^{2}$Department of Medical Imaging and Radiation Sciences, Universit\'e de Sherbrooke, Sherbrooke, QC, Canada}\\
{\small $^{3}$Department of Biochemistry, Faculty of Science, University of Yaound\'e I, P.O. Box 812, Yaound\'e, Cameroon}\par}

\begin{abstract}
Mechanism-of-action (MoA) prediction from molecular structure is often evaluated under random splits that can place chemically identical or closely related compounds on both sides of a test boundary. We constructed an auditable mapping between \num{3289} drug-level LISH-MoA records and \num{1722} unique molecular structures, resolved salts, assigned identical canonical structures to \num{794} collision groups, and evaluated a graph neural network, two graph--descriptor fusion models, and a pretrained molecular transformer under collision-group-disjoint cross-validation. Under this leakage-controlled protocol, the learned molecular arms (GIN, GIN-TFP, GIN-TNE, ChemBERTa) remained close to chance on macro-AUROC (\num{0.501}--\num{0.508}); the fingerprint baseline reached \num{0.536}, and the phenotype-only reference \num{0.63619}. The best learned-arm log loss was obtained by GIN-TFP (\num{0.02334}), but low log loss accompanied near-chance ranking, indicating calibrated but non-discriminative predictions. These results do not demonstrate that molecular representations are universally inferior; they indicate that, under this mapped cohort, endpoint, and evaluation boundary, structure alone did not provide a convincing predictive advantage for MoA-associated labels. The benchmark addresses prediction of observed MoA-associated annotations and does not establish causal mechanism, target engagement, or biological activity.
\end{abstract}

\section{Introduction}

Mechanism-of-action prediction is a multi-label problem in which one compound may be associated with several pharmacological annotations. Molecular structure and cellular phenotype describe different levels of this problem: structure encodes chemical constitution, whereas phenotype summarizes responses observed across cellular systems and experimental conditions \citep{Trapotsi2022MoA}. A useful benchmark should therefore distinguish predictive performance from biological mechanism assignment and should evaluate generalization under boundaries that reflect the intended deployment setting.

Random drug-level cross-validation can be optimistic when repeated records, salts, or identical canonical structures cross train and test sets. Collision-aware evaluation is particularly important for multi-label drug data because labels and covariates can be shared across records that are chemically identical \citep{guo2024scaffold}. Recent benchmarks indicate that representation scale does not guarantee improved molecular prediction: classical machine learning on fingerprints wins \qty{47.4}{\percent} of task--metric entries in a \num{156}-comparison study \citep{guo_ding_2026}, and almost no neural embedding statistically improves over ECFP in a \num{25}-embedding comparison \citep{benchmark_pretrained_2025}. The present study asks whether molecular structure provides transferable predictive information for observed MoA-associated labels beyond a phenotype-only reference when canonical-structure collisions are prevented from crossing the evaluation boundary, and whether combining structure with phenotype improves out-of-sample prediction. We compare a graph model, graph--descriptor fusion models, a pretrained molecular transformer, and a late-concat fusion baseline under a common cohort and fold contract. The analysis is a prediction benchmark, not a causal MoA study.

\section{Materials and methods}

\subsection{Study design and endpoint}

The dataset comprised \num{3289} drug-level rows and \num{206} scored MoA-associated labels from the LISH MoA training release \citep{lish_moa_2019}. The primary endpoint was mean column-wise log loss, computed by averaging binary log losses over labels. Secondary endpoints were macro-AUROC \citep{bmj2006rocauc}, macro-AUPRC \citep{davis2006auprc}, mean Brier score, and mean expected calibration error (ECE) \citep{natoesci2017ece}. AUROC and AUPRC were calculated only for labels with both classes in the test fold; this denominator was retained in the fold records. Probabilities were clipped to $[10^{-7},1-10^{-7}]$ for log-loss evaluation.

\subsection{Structure mapping and collision control}

A versioned mapping from \texttt{drug\_id} to SMILES was constructed from PRISM-aligned records. All \num{3289} drug-level rows were covered. Multi-component entries were split, the largest fragment by heavy-atom count was retained, and structures were canonicalized with RDKit \citep{RDKit}. The validated mapping contained \num{1722} unique molecules and \num{794} collision groups. The frozen mapping contract and its hash are recorded in the project data manifest. Collision groups, rather than individual rows, were assigned to folds; no identical canonical structure was allowed to occur in both training and test data.

\subsection{Models}

Six model configurations were evaluated:

\begin{itemize}
\item \textbf{Phenotype reference}: an unweighted per-label logistic model fitted independently for each of the \num{206} labels, using the predefined gene-expression, viability, and experimental-condition features. This model operates on the same drug-level records but uses no molecular structure information.
\item \textbf{ECFP4--RF}: a random forest on \num{2048}-bit Morgan radius-\num{2} fingerprints, included as a molecular baseline.
\item \textbf{GIN}: a graph isomorphism network \citep{xu2019gin} with \num{128} hidden units, \num{3} layers, ReLU activations, batch normalization, dropout $p=\num{0.1}$, mean-pooling readout, and a multi-label output of \num{206} units.
\item \textbf{GIN--TFP}: the same GIN with tensor-free persistent descriptors (persistent homology, persistent images, and Betti descriptors; \num{78} dimensions) concatenated at the head.
\item \textbf{GIN--TNE}: the same GIN with tensor-network embeddings (\num{192} dimensions) concatenated at the head. Four TNE descriptor failures were retained as zero-vector input rows under an intention-to-treat accounting rule; no failure was removed from the denominator.
\item \textbf{ChemBERTa}: a pretrained SMILES transformer (\texttt{seyonec/ChemBERTa-zinc-base-v1}) \citep{chemberta_2020,chemberta3_2026}, fine-tuned with cross-entropy and independent pretrained-weight initialization for every fold.
\end{itemize}

\subsection{Splits and validation}

The primary evaluation used five seeds and five collision-group-disjoint folds (\num{25} fold--seed replicates). For each molecular fold, model selection used only training groups: training groups were randomly assigned to six parts and one part was used for validation. Early stopping minimized validation mean column-wise log loss with a patience of \num{10} epochs. No test group entered model selection. The phenotype baseline was additionally evaluated under the original drug-level k-fold split and a scaffold-held-out sensitivity split \citep{guo2024scaffold}. All reported molecular-arm results derive from the same mapped cohort and primary fold definition.

\subsection{Training protocol}

All deep models were trained with binary cross-entropy loss over the \num{206} multi-label output units. GIN-family models used AdamW with learning rate \num{1e-3}, weight decay \num{1e-4}, batch size \num{512}, and a maximum of \num{50} epochs. Early stopping monitored validation mean column-wise log loss with a patience of \num{10} epochs; the validation-optimal epoch was selected before test evaluation. ChemBERTa used AdamW with learning rate $2\times 10^{-5}$, weight decay \num{0.01}, batch size \num{32}, maximum sequence length \num{128}, and a maximum of \num{10} epochs with patience \num{3}. For each fold, the pretrained ChemBERTa state was restored before fine-tuning, ensuring no cross-fold weight transfer. All deep models ran on a single NVIDIA RTX A4000 GPU. The ECFP4--RF baseline used scikit-learn's \texttt{RandomForestClassifier} with \num{500} trees and default hyperparameters. The phenotype logistic baseline was fitted independently for each label with L2 regularization.

\subsection{Reproducibility and analysis boundaries}

Mapping records, descriptor computations, fold assignments, and per-fold model outputs were preserved with their random seeds, collision policy, and descriptor-failure accounting. Two extensions to the primary protocol are reported as bounded proxies under the same scaffold split: pooled calibration metrics and a quantile-kernel-similarity (QKS) separability diagnostic. The fold-level prediction files used as input are gated on a documented column contract and on MoA-set invariance across folds; results that fail this audit are not aggregated. The present calibration analysis is pooled across labels; per-label reliability curves and calibration slope/intercept are outside the current estimand. The QKS result is likewise a bounded proxy rather than a full quantum-kernel implementation.

\section{Results}

\subsection{Mapping audit}

The mapping covered all \num{3289} drug-level rows. After salt stripping and canonicalization, all structures passed the RDKit validity gate. The final cohort contained \num{1722} unique molecules and \num{794} collision groups. These collision groups formed the unit of assignment for the primary evaluation.

\subsection{Primary collision-group benchmark}

\Cref{tab:p6_primary} presents the collision-group results across \num{25} seed--fold evaluations. GIN-TFP had the lowest learned-arm log loss at \num{0.02334}, followed by GIN at \num{0.02419} and GIN-TNE at \num{0.02434}. ChemBERTa had the highest learned-arm log loss at \num{0.02549}. Macro-AUROC values for the learned molecular arms ranged from \num{0.50137} to \num{0.50785}, the fingerprint baseline reached \num{0.53562}, and the phenotype reference \num{0.63619}. The learned-arm AUPRC values were also substantially lower than the phenotype reference.

\begin{table}[htbp]
\centering
\caption{Audited collision-group results across \num{25} seed--fold evaluations. Log loss is the primary endpoint.}
\label{tab:p6_primary}
\begin{tabular}{lrrrrr}
\toprule
Arm & Log loss & Macro-AUROC & Macro-AUPRC & Brier & ECE\\
\midrule
Phenotype reference & \num{0.02428} & \num{0.63619} & \num{0.13137} & \num{0.00381} & \num{0.00387}\\
ECFP4--RF & \num{0.03582} & \num{0.53562} & \num{0.02590} & \num{0.00344} & \num{0.00232}\\
GIN & \num{0.02419} & \num{0.50785} & \num{0.01379} & \num{0.00347} & \num{0.00449}\\
GIN-TFP & \num{0.02334} & \num{0.50479} & \num{0.01381} & \num{0.00341} & \num{0.00300}\\
GIN-TNE & \num{0.02434} & \num{0.50601} & \num{0.01460} & \num{0.00341} & \num{0.00306}\\
ChemBERTa & \num{0.02549} & \num{0.50137} & \num{0.01313} & \num{0.00343} & \num{0.00750}\\
\bottomrule
\end{tabular}
\end{table}

The dissociation between log loss and AUROC is informative: a low multi-label log loss can be achieved by well-calibrated prevalence-like predictions, whereas ranking ability requires discriminating which drugs carry which labels. The structure arms therefore appear reasonably calibrated but non-discriminative on this task. Paired seed--fold bootstrap intervals for the scaffold comparisons are provided in the data release; they quantify variation across the 25 predefined replicates and are not population-level confidence intervals over molecules. These comparisons are exploratory and are not used to support a superiority claim.

\subsection{Reference and sensitivity analyses}

\Cref{tab:p6_splits} compares performance across the three evaluated splits for the phenotype reference and ECFP4--RF baseline. The phenotype reference reproduced the locked random-split result (AUROC \num{0.64345}) and showed minimal degradation under collision-group splitting (\num{0.63619}) and scaffold sensitivity (\num{0.64023}). The ECFP4--RF baseline reached AUROC \num{0.53562} under collision-group splitting and \num{0.53817} under scaffold splitting; it was not evaluated under the drug-level random split. The phenotype--ECFP4 fusion arm reached AUROC \num{0.58323}, below the phenotype reference.

\begin{table}[htbp]
\centering
\caption{Split sensitivity for the phenotype reference and ECFP4--RF baseline.}
\label{tab:p6_splits}
\begin{tabular}{lccc}
\toprule
Arm & Random & Collision-group & Scaffold \\
\midrule
Phenotype reference & \num{0.64345} & \num{0.63619} & \num{0.64023}\\
ECFP4--RF & --- & \num{0.53562} & \num{0.53817}\\
\bottomrule
\end{tabular}
\end{table}

These comparisons reinforce the importance of reporting the split and endpoint together. The modest random-to-collision-group cost for the phenotype reference (\num{0.00726} AUROC units) reflects the limited impact of collision awareness on this phenotype-driven task, whereas the structure baselines remain near chance across all splits.

\subsection{Scaffold-held-out molecular arms}

The four molecular arms were also evaluated under a stricter scaffold-held-out boundary, in which whole Bemis--Murcko scaffolds were kept out of training. Mean metrics over the same \num{25} fold--seed grid are reported in \Cref{tab:p6_scaffold_gnn}. All four arms remained at chance macro-AUROC (\num{0.501}--\num{0.506}), below the scaffold phenotype reference (\num{0.64023}) and reproducing the collision-group ordering. The best scaffold molecular-arm log loss was again obtained by GIN-TFP (\num{0.02373}), marginally below its collision-group value, while its macro-AUROC (\num{0.50077}) stayed at chance. The near-chance ranking therefore extends to a stricter chemical-generalisation boundary and is not an artefact of the collision-group split alone.

\begin{table}[htbp]
\centering
\caption{Scaffold-held-out results for the four molecular arms across \num{25} seed--fold evaluations. Log loss is the primary endpoint.}
\label{tab:p6_scaffold_gnn}
\begin{tabular}{lrrrrr}
\toprule
Arm & Log loss & Macro-AUROC & Macro-AUPRC & Brier & ECE\\
\midrule
GIN & \num{0.02474} & \num{0.50642} & \num{0.01509} & \num{0.00349} & \num{0.00493}\\
GIN-TFP & \num{0.02373} & \num{0.50077} & \num{0.01300} & \num{0.00343} & \num{0.00325}\\
GIN-TNE & \num{0.02429} & \num{0.50609} & \num{0.01399} & \num{0.00342} & \num{0.00318}\\
ChemBERTa & \num{0.02559} & \num{0.50344} & \num{0.01327} & \num{0.00344} & \num{0.00761}\\
\bottomrule
\end{tabular}
\end{table}

The log-loss/AUROC dissociation observed under collision-group folds therefore persists under scaffold-held-out extraction: the molecular arms produce well-calibrated, prevalence-like probabilities that approach the marginal label distribution without ranking drugs of different MoA-associated profiles.

\subsection{Calibration diagnostics}

\Cref{tab:p6_calibration} reports pooled calibration metrics computed over the archived scaffold-split fold-level test predictions (\num{25} seed--fold replicates, $n=\num{3387670}$ drug--label observations pooled across \num{206} MoA labels). The analysis used only fold-level prediction files matching the documented \texttt{drug\_id}, \texttt{<moa>\_true}, \texttt{<moa>\_pred} column contract and seed/fold convention, with an invariant MoA set across folds. The phenotype and structure arms are tightly calibrated (pooled ECE \num{0.0016} and \num{0.0012}; Brier \num{0.0038} and \num{0.0040}), and the fusion arm remains well calibrated (ECE \num{0.0017}; Brier \num{0.0034}). The maximum calibration error is dominated by a single tail bin in every arm, indicating that the pooled ECE is not masking a localized miscalibration. The archived JSON files also contain per-label Brier summaries and the prevalence-tertile sensitivity stratification described in the DAR; those records are used for auditability, not as a claim of label-specific calibration. These results are consistent with the collision-group conclusion that the molecular arms produce prevalence-like, well-calibrated probabilities rather than discriminative rankings.

\begin{table}[htbp]
\centering
\caption{Pooled calibration metrics over scaffold-split fold-level test predictions (\num{25} replicates, $n=\num{3387670}$ drug--label pairs, \num{206} MoA labels).}
\label{tab:p6_calibration}
\begin{tabular}{lrrrr}
\toprule
Arm & $n$ & ECE & MCE & Brier \\
\midrule
Phenotype reference & \num{3387670} & \num{0.0016} & \num{0.648} & \num{0.0038} \\
ECFP4--RF (structure) & \num{3387670} & \num{0.0012} & \num{0.389} & \num{0.0034} \\
Phenotype\,+\,ECFP4--RF fusion & \num{3387670} & \num{0.0017} & \num{0.259} & \num{0.0034} \\
\bottomrule
\end{tabular}
\end{table}

\subsection{Bounded QKS separability}

\Cref{tab:p6_qks} reports a bounded quantile-kernel-similarity (QKS) proxy between the phenotype arm and the two molecular arms over the same \num{3289} scaffold-split drugs. The proxy is the absolute difference of empirical CDF quantiles of the per-drug maximum predicted probability, evaluated at $\tau\in\{0.25,0.50,0.75\}$, supplemented with the Spearman rank correlation between cross-arm maximum predicted probability. The protocol is bounded to a single drug cohort and to the scaffold split; it is not a reference quantum-kernel implementation. Values are taken from the extended four-pair run archived with the analysis artifacts. The phenotype-vs-structure proxy at $\tau=0.5$ is $\num{0.023}$ and the cross-arm Spearman on maximum predicted probability is $\num{0.036}$, indicating that the structure arm does not add ranking information on top of the phenotype arm. The phenotype-vs-fusion proxy at $\tau=0.5$ is $\num{0.033}$ and the cross-arm Spearman is $\num{0.427}$, indicating moderate agreement between the phenotype arm and the fusion arm on the predicted-probability ranking; this descriptive association does not establish incremental predictive value.

\begin{table}[htbp]
\centering
\caption{Bounded QKS proxy between the phenotype arm and the two molecular arms over the scaffold-split prediction set (\num{3289} drugs, $\tau=0.5$ headline; $\Delta F(0.5)=\lvert F_p(0.5)-F_y(0.5)\rvert$, the absolute difference of empirical CDF quantiles of per-drug maximum predicted probability).}
\label{tab:p6_qks}
\begin{tabular}{lrr}
\toprule
Pair & $\Delta F(0.5)$ & $\rho_S$ max-prob \\
\midrule
Phenotype vs.\ structure (ECFP4)   & \num{0.023} & \num{0.036} \\
Phenotype vs.\ fusion (phenotype $+$ ECFP4) & \num{0.033} & \num{0.427} \\
\bottomrule
\end{tabular}
\end{table}

\section{Discussion}

This benchmark was designed to test whether molecular structure provides transferable information for predicting observed MoA-associated labels beyond a phenotype-only reference, and whether a simple combination of the two sources improves prediction under leakage-controlled evaluation. For the evaluated molecular arms, structure did not recover the phenotype ranking signal: the negative result persists under both collision-group (\Cref{tab:p6_primary}) and scaffold-held-out (\Cref{tab:p6_scaffold_gnn}) splits, while the late-concat fusion remains below the phenotype reference.

\subsection{Calibrated but non-discriminative}

GIN-TFP achieved the lowest molecular-arm log loss (\num{0.02334} under collision group; \num{0.02373} under scaffold), marginally below even the phenotype reference on this endpoint, but its macro-AUROC (\num{0.50479} under collision group; \num{0.50077} under scaffold) remained close to chance and far from the phenotype reference (\num{0.63619}; \num{0.64023} scaffold). This dissociation is informative: a low multi-label log loss can be achieved by well-calibrated prevalence-like predictions, whereas ranking ability requires discriminating which drugs carry which labels. The structure arms therefore appear reasonably calibrated but non-discriminative on this task under both evaluation boundaries.

\subsection{Why phenotype features dominate}

The phenotype reference performed better on the ranking metrics, which is expected because the labels are associated with cellular response annotations and the phenotype features summarize cellular measurements. This does not make the phenotype model a causal mechanism predictor. Conversely, the structure models do not fail because they are disproven in general; they may require larger chemically diverse cohorts, improved molecular supervision, alternative label definitions, or task-specific pretraining. The present result is therefore an honest benchmark conclusion rather than a universal model ranking.

\subsection{Comparison with P5 molecular benchmark}

The same model architectures were evaluated on a different task in a companion study of antimalarial activity prediction \citep{temgoua2027quantum}. Under scaffold-held-out evaluation on \num{19836} African antimalarial natural products, ECFP4--RF achieved ROC AUC \num{0.8300}, exceeding all learned arms. The present P6 benchmark differs in task (MoA labels vs.\ binary antimalarial activity), label source (phenotype annotations vs.\ screening activity), and cohort composition (drug-level LISH records vs.\ natural-product screening data). The consistent ECFP4--RF advantage across both tasks supports the interpretation that local substructure encodings captured by fingerprints remain highly data-efficient for small-to-medium molecular panels, regardless of whether the endpoint is activity classification or MoA-label prediction.

\subsection{Implications for representation choice}

For comparable drug-level MoA benchmarks, the tested molecular representations did not provide a ranking advantage over the phenotype reference. This result is consistent with broader evidence that classical fingerprint models often remain competitive with graph and sequence models on small-to-medium datasets \citep{guo_ding_2026,benchmark_pretrained_2025}. The implication is practical: practitioners evaluating molecular representations for MoA-associated label prediction should test against a phenotype baseline under leakage-controlled splits before assuming that more complex representations will improve performance.

The dissociation between log loss and AUROC across all molecular arms is notable. Low log loss with near-chance AUROC indicates that the models produce calibrated probability estimates that approximate the marginal label distribution without learning to discriminate among drugs. This pattern suggests that the molecular features, as currently featurized, do not carry sufficient task-specific signal to override the prevalence-driven baseline for this multi-label phenotype task.

\subsection{Collision groups as a contribution}

The collision-group mapping is a central contribution of the analysis. It prevents identical canonical structures from appearing in both training and test folds, while retaining duplicated drug identities as an explicit part of the data-generating record. The mapping also exposes a limitation: structure collisions are common, and the effective chemical diversity is closer to \num{1722} unique molecules than to the raw row count of \num{3289}. Four TNE failures were retained as zero-vector intention-to-treat cases, preserving the declared denominator but potentially attenuating descriptor-specific performance.

\section{Limitations}

First, the benchmark predicts MoA-associated annotations rather than experimentally established causal mechanisms. No target engagement, antimalarial activity, resistance phenotype, or prospective biological validation was measured. Second, the primary molecular comparison used collision-group folds; a scaffold-held-out sensitivity is reported for the four GNN and ChemBERTa arms and preserves the near-chance verdict, but it remains a sensitivity analysis on the same mapped cohort rather than a prospectively independent chemical-generalisation test. Third, the paired seed--fold analysis quantifies variation across the 25 predefined replicates, but it is not a substitute for independent external validation. The comparisons are exploratory; no confirmatory multiplicity claim or population-level confidence interval is made. Fourth, the archived outputs contain per-label Brier summaries, but reliability diagrams and calibration slope/intercept were not computed; the pooled ECE, MCE, and Brier score reported in the Results section are summary calibration diagnostics only, and the prevalence-tertile stratification is a sensitivity analysis rather than a biological label classification. Fifth, the bounded QKS protocol is a CDF-quantile proxy, not a full quantum-kernel reference implementation, and is reported as a single-fingerprint separability diagnostic. Sixth, the mapping relies on PRISM-aligned source records and a frozen matching procedure; external structure-source replication may alter coverage or collision composition, and only \num{100}{\%} signature-matched rows were included.

\section{Future directions}

Several extensions could strengthen the evidence base. First, a truly prospective chemical-generalisation test---for example scaffold splits on an independently mapped cohort or experimental validation of the highest-confidence predictions---would test whether the near-chance pattern persists beyond the frozen PRISM-aligned cohort; the current scaffold sensitivity reuses the same cohort and mapping. Second, reliability diagrams, calibration slope/intercept, and per-MoA ECE could be added in a dedicated calibration study; the present work retains per-label Brier summaries and a prevalence-tertile sensitivity analysis but does not interpret either as evidence of label-specific biological calibration. Third, contrastive pretraining on the molecular panel could improve representation quality before fine-tuning. Fourth, the bounded QKS proxy could be replaced with a full reference kernel implementation once the per-drug descriptor matrix is finalized. Finally, prospective biochemical validation on a subset of high-confidence predictions would test whether the computational ranking predicts biological activity.

\section{Conclusion}

Under a frozen, leakage-controlled evaluation of \num{3289} LISH-MoA drug-level records, the evaluated molecular representations did not provide transferable ranking information beyond the phenotype reference, and the late-concat fusion did not improve on that reference. GIN-TFP obtained the best molecular-arm log loss, but molecular AUROCs remained close to chance. The result supports evaluating structure and phenotype at their appropriate information levels, with collision control and explicit failure accounting; it does not establish universal molecular-model inferiority or causal biological conclusions.

\section*{Data and code availability}

All code, frozen splits, per-fold results, and the analysis code supporting the conclusions of this study are available under the MIT license in the public GitHub repository (\url{https://github.com/NanaEngo/Malaria_codesV2}). A versioned snapshot will be archived on Zenodo to provide a permanent, citable DOI concurrent with publication.

\section*{Author contributions}

Myke Vital Sao Temgoua: Methodology, software, formal analysis, investigation, data curation, writing -- original draft. Jean-Pierre Tchapet Njafa: Methodology, software, supervision, writing -- review \& editing. Penabei Samafou: data curation, writing -- review \& editing. Wilfred Fon Mbacham: conceptualization, validation, writing -- review \& editing. Serge Guy Nana Engo: conceptualization, methodology, supervision, writing -- review \& editing. All authors read and approved the final manuscript.

\section*{Competing interests}

The authors declare no competing interests.

\section*{Funding}

This research received no specific grant from any funding agency in the public, commercial, or not-for-profit sectors.

\section*{Ethics approval and consent to participate}

Not applicable. This study is entirely computational and uses publicly available chemical databases; no human or animal subjects were involved.

\section*{Data and code accuracy statement}

The numerical values in the results tables were recomputed from the released per-fold prediction files and corresponding summary artifacts. Frozen splits, per-fold outputs, descriptors, and prediction files are released with the repository.

\section*{Use of Artificial Intelligence}

The authors used AI-assisted tools during code development and data-analysis scripting. All affected text and code were reviewed and edited by the authors, who take full responsibility for the article. No AI system contributed as an author.

\section*{Acknowledgements}

Computational resources were provided by the University of Yaound\'e I HPC facility. We acknowledge the LISH MoA challenge organizers for the public dataset.

\bibliographystyle{plainnat}
\bibliography{Bibliography_P6}
\end{document}
